% ---------------------------------------------------------------------------
%  QML Lab research project 07: Quantum GAN for a One-Dimensional Distribution
%  Build:  latexmk -pdf main.tex   (needs qmlproject.sty, harvardqml.bst, references.bib)
% ---------------------------------------------------------------------------
\documentclass[11pt,a4paper]{article}
\usepackage{qmlproject}
\qmlheader{Project 07: Quantum GAN}

\hypersetup{pdftitle={QML Lab project 07: Quantum GAN for a One-Dimensional Distribution}}

\begin{document}

\begin{center}
  {\small\textcolor{black!55}{QML Lab \textbullet\ research project 07}}\\[0.5em]
  {\LARGE\bfseries Quantum GAN for a One-Dimensional Distribution}\\[0.6em]
  {\large Learning and loading a probability distribution with a quantum generator.}\\[0.6em]
  {\small Living document \textbullet\ last updated \today}
\end{center}

\begin{projectcard}
\small
\begin{tabularx}{\linewidth}{@{}>{\bfseries}l X >{\bfseries}l X@{}}
Quadrant  & \CQ\ (classical data, quantum processor) & Track & Variational \\
Difficulty & Hard & Tools & Qiskit Machine Learning, PyTorch \\
Builds on & Modules 5 and 7 of the QML Lab & Workspace & \path{qml_projects/quantum-gan-1d/} \\
\end{tabularx}
\end{projectcard}

\paragraph{Summary.}
A quantum generative adversarial network (qGAN) trains a parameterised circuit whose measurement
outcomes follow a target distribution, guided by a classical discriminator. Besides generating
samples, the trained circuit loads the distribution into amplitudes, which other quantum algorithms
need as input. The project studies how well a qGAN learns discretised one-dimensional distributions,
how stable the training is, and how it compares with small classical generators.

% ===========================================================================
\section{Background and key concepts}
% ===========================================================================
\subsection{Generative adversarial networks}
A GAN \citep{goodfellow2014} trains a generator $G$ and a discriminator $D$ in a two-player game,
\begin{equation}
  \min_G\max_D\ \E_{x\sim p_{\text{data}}}\big[\log D(x)\big]+\E_{z}\big[\log\big(1-D(G(z))\big)\big].
\end{equation}
The generator improves until the discriminator can no longer tell its samples from real data.

\subsection{Quantum generators}
Quantum versions were proposed by \citet{lloyd2018qgan} and \citet{dallairedemers2018}. In the qGAN of
\citet{zoufal2019} the generator is an $n$-qubit circuit $U(\theta)$ whose output state
\begin{equation}
  \ket{g_\theta}=U(\theta)\ket{0}^{\otimes n}=\sum_{j=0}^{2^n-1}\sqrt{p_\theta(j)}\,e^{i\varphi_j}\ket{j}
\end{equation}
defines a distribution over $2^n$ bins: measuring returns bin $j$ with probability $p_\theta(j)$. The
discriminator is a classical neural network that sees the sampled bin values. After training, the
circuit is an efficient loader for the learned distribution.

\subsection{Measuring success}
Quality is measured by the relative entropy (Kullback--Leibler divergence) to the target $q$,
\begin{equation}
  D_{\mathrm{KL}}(q\,\Vert\,p_\theta)=\sum_jq(j)\log\frac{q(j)}{p_\theta(j)},
\end{equation}
and by the total variation distance. Training curves often oscillate, as in classical GANs.

% ===========================================================================
\section{Research questions}
% ===========================================================================
The project is open-ended: these questions are the starting point, not the finish line. Add new
ones to the README and the journal as they come up.

\begin{rqlist}
  \item How many qubits (bins) are needed for a useful approximation of log-normal and bimodal targets?
  \item Is training stable, or does it oscillate or collapse to one mode?
  \item How does the KL divergence compare with a small classical GAN and with a directly parameterised histogram?
\end{rqlist}

% ===========================================================================
\section{Method}
% ===========================================================================
\subsection{Data}
Targets discretised to $2^n$ bins: a log-normal distribution (as in \citealp{zoufal2019}), a
truncated normal and a bimodal mixture. Training samples are drawn from the discretised targets.

\subsection{Models}
Generator: hardware-efficient circuit with $R_Y$ rotations and CZ entanglers, $L$ layers, initialised
from a uniform or normal-like distribution. Discriminator: small PyTorch MLP \citep{paszke2019}.
Alternating updates with Adam; generator gradients through Qiskit Machine Learning's PyTorch connector.

\subsection{Baselines}
A classical GAN with a small MLP generator; a softmax-parameterised histogram with $2^n-1$
parameters fitted by maximum likelihood (the best possible fit for the bins).

\subsection{Metrics}
\begin{itemize}
  \item KL divergence and total variation distance to the target over training.
  \item Oscillation amplitude of the loss curves (training stability).
  \item Number of parameters and circuit evaluations.
\end{itemize}

% ===========================================================================
\section{Planned experiments}
% ===========================================================================
Each experiment becomes one \texttt{[[experiments]]} entry in the workspace, so its runs, parameters
and metrics are stored and compared in the portal.

\begin{table}[h]
\centering\small
\begin{tabularx}{\linewidth}{@{}l>{\raggedright}X>{\raggedright}X>{\raggedright\arraybackslash}X@{}}
\toprule
ID & Experiment & What varies & Main metrics \\
\midrule
  E1 & Log-normal, 3 qubits & epochs & KL divergence \\
  E2 & Resolution & $n=3,4,5$ qubits & KL, cost \\
  E3 & Bimodal target & mixture weights & mode coverage \\
  E4 & Generator depth & $L=1..4$ & KL, stability \\
  E5 & Shots per step & $M\in\{100,1000\}$ & KL, stability \\
\bottomrule
\end{tabularx}
\caption{Experiment plan. Results go into the journal and the results log (\cref{sec:log}).}
\end{table}

% ===========================================================================
\section{Running it in the QML Lab}
% ===========================================================================
The workspace lives in \path{qml_projects/quantum-gan-1d/}. Scripts import three helpers from
\path{qml_lab/plots.py}:
\begin{itemize}
  \item \texttt{params()} reads the parameters from the portal's run form (defaults when run by hand);
  \item \texttt{save()} stores figures with the run;
  \item \texttt{log\_metrics()} records numbers that appear in the run table.
\end{itemize}

\begin{lstlisting}
qml_projects/quantum-gan-1d/
  project.toml          metadata + experiments
  README.md             goal, open questions, ideas
  experiments/          runnable scripts
  docs/main.tex         this document
\end{lstlisting}

First experiment: This experiment is \emph{planned}: write the script, add the block below to \texttt{project.toml}, then run \texttt{python manage.py sync\_projects}.

\begin{lstlisting}[language={}]
[[experiments]]
slug = "lognormal"
title = "qGAN for a 3-qubit log-normal distribution"
script = "experiments/e01_lognormal.py"
figure = "distribution.png"

[experiments.params]
n_qubits = 3
n_layers = 2
epochs = 300
shots = 1000
target = "lognormal"
seed = 0
\end{lstlisting}

\begin{lstlisting}[language=bash]
# from the portal folder: pick up changes to project.toml
python manage.py sync_projects
cd qml_projects/quantum-gan-1d
# run by hand with the default parameters
PYTHONPATH=../../qml_lab python experiments/e01_lognormal.py
\end{lstlisting}

% ===========================================================================
\section{Quality bar}
% ===========================================================================
\begin{itemize}
  \item A classical baseline with a similar number of trainable parameters, tuned with the same care \citep{bowles2024}.
  \item At least five random seeds per configuration, reported as mean $\pm$ standard deviation.
  \item How results change with the number of qubits.
  \item At least one run under a realistic noise model.
  \item Show the learned histogram next to the target and the classical fits in every report.
\end{itemize}

% ===========================================================================
\section{Open directions}
% ===========================================================================
\begin{itemize}
  \item Use the trained loader as input to amplitude estimation, the original application of \citet{zoufal2019}.
  \item Quantum circuit Born machines trained with maximum mean discrepancy instead of a discriminator.
  \item Two-dimensional distributions on more qubits.
\end{itemize}

% ===========================================================================
\section{Risks and limitations}
% ===========================================================================
\begin{itemize}
  \item GAN training is unstable; use several seeds and report the spread.
  \item For 8 bins a classical histogram fits exactly; the interest is in loading and scaling, not in beating it.
\end{itemize}

% ===========================================================================
\section{Results log}
\label{sec:log}
% ===========================================================================
Summarise findings here as they become stable; day-to-day notes belong in the portal journal.

\begin{table}[h]
\centering\small
\begin{tabularx}{\linewidth}{@{}l l >{\raggedright}X >{\raggedright\arraybackslash}X@{}}
\toprule
Date & Run & Finding & Next step \\
\midrule
 & & & \\
 & & & \\
 & & & \\
\bottomrule
\end{tabularx}
\end{table}

\bibliographystyle{harvardqml}
\bibliography{references}

\end{document}
